\chapter{Knowledge Editing Methodologies}

% \section{Steering Vectors?}

% TODO

\section{Soft Prompting}
\label{sec:soft_prompting}

Soft prompting optimises a small set of continuous vectors, prepended to the
input, while keeping all base model parameters frozen. The motivation is to test
whether a low-dimensional perturbation at the input is sufficient to shift the
model's residual stream into a state that supports better arithmetic computation,
and whether the resulting shift passes through the same circuit identified by the
contrastive analysis.

\subsection{Formulation}
\label{sec:soft_prompt_formulation}

Let $x = (t_1, \ldots, t_T)$ be a prompt tokenised to a sequence of length $T$,
with embedding matrix $E \in \mathbb{R}^{V \times d}$ mapping each token to a
$d$-dimensional vector. Two variants are implemented, differing in where the
learned vectors are injected.

\paragraph{Soft prompt.}
A parameter matrix $P \in \mathbb{R}^{k \times d}$ is introduced, where $k$ is
the prefix length. At the embedding layer, the input sequence is extended to
\begin{equation}
\bigl(P_1, \ldots, P_k,\; e_{t_1}, \ldots, e_{t_T}\bigr) \in \mathbb{R}^{(k+T) \times d},
\label{eq:soft_prompt_seq}
\end{equation}
where $e_{t_i} = E[t_i]$ are the real token embeddings. All subsequent layers
process this extended sequence; the prefix vectors attend to and are attended to
by all real tokens. The only trainable parameters are the $k \times d$ entries of
$P$.

\paragraph{Prefix tuning.}
A separate parameter matrix $P^{(l)} \in \mathbb{R}^{k \times d}$ is introduced
per layer $l \in \{0, \ldots, L-1\}$, collected as $\mathcal{P} \in
\mathbb{R}^{L \times k \times d}$. Before each layer's self-attention, the first
$k$ positions of the residual stream are overwritten with $P^{(l)}$:
\begin{equation}
h^{(l)}_{1:k} \leftarrow P^{(l)}.
\label{eq:prefix_tuning_inject}
\end{equation}
This forces the prefix positions to hold fixed, learned values at every layer
rather than allowing them to evolve through the network. The total parameter count
is $L \times k \times d$, compared with $k \times d$ for the soft prompt.

\paragraph{Implementation.}
Both variants are implemented without modifying the base model's forward pass.
Placeholder pad tokens are prepended to the real token IDs to extend the sequence
to length $k + T$, and the attention mask is extended with $k$ ones so that
prefix positions are attended to normally. A forward hook then overwrites the
first $k$ positions at the appropriate point, either at the embedding
output for the soft prompt or at each layer's residual stream input for prefix
tuning.

\subsection{Training}
\label{sec:soft_prompt_training}

All base model parameters are frozen. Only $P$ (or $\mathcal{P}$) receives
gradients. The training objective is cross-entropy on the first answer token at
the last prompt position. Given a batch of prompts with extended inputs
$(P \oplus e_{t_1:t_T})$ and target token indices $y$, the loss is
\begin{equation}
\mathcal{L} = -\frac{1}{B}\sum_{b=1}^{B} \log \frac{\exp f_{y_b}(x_b)}{\sum_{v} \exp f_v(x_b)},
\label{eq:soft_prompt_loss}
\end{equation}
where $f_v(x_b)$ is the logit for token $v$ at the final prompt position.
The optimiser is AdamW with weight decay $10^{-4}$ and learning rate $10^{-3}$,
with a cosine annealing schedule over the training epochs. Gradients are clipped
to unit norm. The dataset is split 90/10 into training and validation sets; the
checkpoint with the lowest validation cross-entropy is retained.

\subsection{Alignment Analysis}
\label{sec:soft_prompt_alignment}

After training, the experiment tests whether the prefix improves accuracy by
amplifying the same concept direction identified by the contrastive localisation
analysis. For each layer $l$ in a set of sweep layers, the \emph{prefix delta} is
computed as the mean difference between the residual stream at the answer-token
position with and without the prefix,
\begin{equation}
\Delta^{(l)}
=
\frac{1}{N} \sum_{i=1}^{N}
\Bigl[
  h^{(l)}_{\mathrm{pfx}}(x_i)
  -
  h^{(l)}_{\mathrm{base}}(x_i)
\Bigr],
\label{eq:prefix_delta}
\end{equation}
where $h^{(l)}_{\mathrm{pfx}}$ and $h^{(l)}_{\mathrm{base}}$ are the
post-block residual streams at the answer position, with and without the prefix
hooks active. Two contrastive steering vectors are also computed at the same layers.
The first, $\mathrm{sv}_{\mathrm{base}}^{(l)} = \bar{h}^{(l)}_{\mathrm{correct}} -
\bar{h}^{(l)}_{\mathrm{wrong}}$, is averaged over baseline runs bucketed by whether
the baseline prediction was correct; the second, $\mathrm{sv}_{\mathrm{pfx}}^{(l)}$,
is computed analogously from prefix runs. The cosine similarities
$\cos(\Delta^{(l)}, \mathrm{sv}_{\mathrm{base}}^{(l)})$ and
$\cos(\Delta^{(l)}, \mathrm{sv}_{\mathrm{pfx}}^{(l)})$ per layer indicate
whether the prefix shifts the residual stream in the same direction as the
correct-minus-wrong contrastive vector. Finally, $\Delta^{(l)}$ is encoded
through the layer-$l$ transcoder to identify which sparse features carry the
largest component of the prefix's effect, and these are cross-referenced against
the concept-localisation features found in Chapter~\ref{sec:concept_localisation}.

\subsection{Results}
\label{sec:soft_prompt_results}

The soft prompt is evaluated on three concepts, each representing a different type
of required computation. These are residue class membership (a modular arithmetic
relation), balanced parentheses (a recursive structural property), and causal
direction (a semantic plausibility judgement). For each concept the training prompt template
is T0, the prefix length is $k=10$, and the model is Qwen3-4B with all base parameters
frozen. Accuracy and mean $p(\text{correct})$ are reported on the held-out 10\%
validation split. Table~\ref{tab:sp_summary} summarises the results across all three concepts.

\begin{table}[ht!]
\centering
\small
\begin{tabular}{@{}lrrrr@{}}
\toprule
Concept & Base acc.\ (\%) & Prefix acc.\ (\%) & $\Delta$ acc.\ (\%) & $\Delta\,\bar{p}(\text{correct})$ \\
\midrule
Residue class        & 15.0 & 16.0 & $+1.0$  & $-0.011$ \\
Balanced parentheses &  0.0 & 28.6 & $+28.6$ & $+0.078$ \\
Causal direction     &  0.0 & 53.8 & $+53.8$ & $+0.236$ \\
\bottomrule
\end{tabular}
\caption{Soft prompt results across three concepts. Template T0, prefix length $k=10$,
held-out 10\% validation split ($n=100$, $7$, $13$ respectively).
$\Delta\,\bar{p}$ is the change in mean softmax probability assigned to the correct
first answer token.}
\label{tab:sp_summary}
\end{table}

\textbf{Residue Class (\texttt{calc: \{a\}\%7=})}

The residue class prefix produces no meaningful improvement. First-token accuracy
increases from 15\% to 16\% ($+1$ percentage point), and the mean probability
assigned to the correct token \emph{decreases} by $0.011$.
The baseline model applies a strong surface heuristic, predicting the leading digit
of the operand as the answer, which is correct when $a$ begins with its
own residue but otherwise fails. The prefix cannot overcome this bias. On most prompts it reduces confidence
across all tokens without redirecting probability mass toward the correct residue.
The null result here is consistent with the concept localisation findings
(Section~\ref{sec:residue_results}): no anchor position produces a delta norm
that exceeds the null permutation band, suggesting that the model does not form a
clean binary representation of $a\bmod 7$ at any single token position or depth.

The following extracts illustrate representative behaviour on the validation split.
The predicted token and its softmax probability are shown for the baseline and the
prefix condition; the correct answer is in square brackets.

\begin{quote}
\fontfamily{pcr}\selectfont\small
\textbf{Prompt:} \textup{calc: 134\%7= }\hfill[\textit{correct: }1]\\
\textbf{Baseline:}\ `1' (0.68) \quad `3' (0.12) \quad `2' (0.08)\\
\textbf{Prefix:}\ \ \ `1' (0.32) \quad `3' (0.21) \quad `2' (0.15)\\
\textit{Prediction unchanged; confidence reduced.}\\[6pt]
%
\textbf{Prompt:} \textup{calc: 925\%7= }\hfill[\textit{correct: }1]\\
\textbf{Baseline:}\ `9' (0.64) \quad `1' (0.04) \quad `2' (0.08)\\
\textbf{Prefix:}\ \ \ `1' (0.13) \quad `9' (0.12) \quad `6' (0.10)\\
\textit{Prefix corrects the prediction ($p$ rises from 0.036 to 0.130).}\\[6pt]
%
\textbf{Prompt:} \textup{calc: 803\%7= }\hfill[\textit{correct: }5]\\
\textbf{Baseline:}\ `8' (0.67) \quad `5' (0.01) \quad `3' (0.10)\\
\textbf{Prefix:}\ \ \ `8' (0.16) \quad `5' (0.06) \quad `3' (0.09)\\
\textit{Both conditions wrong; prefix reduces leading-digit bias without correcting answer.}
\end{quote}

\textbf{Balanced Parentheses (\texttt{Is \{seq\} a valid bracket string? Answer yes or no:})}

Accuracy rises from 0\% to 28.6\% (2 of 7 validation prompts) and mean $p(\text{correct})$
increases by $+0.078$.
The baseline model entirely fails to produce the required output format. Its top-1
prediction is the token \texttt{`1'} on nearly every prompt, an artefact of the
model not recognising the yes/no response format from this template alone.
The prefix shifts the output distribution toward the tokens \texttt{`yes'} and
\texttt{`no'}, and in 2 of 7 cases selects the correct answer.
The validation set is small ($n=7$), so the absolute accuracy numbers carry
limited statistical weight; the increase in $\bar{p}(\text{correct})$ provides
a more stable signal.

\begin{quote}
\fontfamily{pcr}\selectfont\small
\textbf{Prompt:} \textup{Is ( ( ) ( ( ) ) ) a valid bracket string? Answer yes or no: }\hfill[\textit{correct: yes}]\\
\textbf{Baseline:}\ `1' (0.25) \quad `Yes' (0.19) \quad `No' (0.17)\\
\textbf{Prefix:}\ \ \ `No' (0.25) \quad `Yes' (0.18) \quad \ldots\\
\textit{Both conditions wrong; prefix formats output correctly but picks wrong label.}\\[6pt]
%
\textbf{Prompt:} \textup{Is ( ( ( ( ) ) ( ) a valid bracket string? Answer yes or no: }\hfill[\textit{correct: no}]\\
\textbf{Baseline:}\ `1' (0.24) \quad `No' (0.24) \quad \ldots\\
\textbf{Prefix:}\ \ \ `No' (0.32) \quad `Yes' (0.14) \quad \ldots\\
\textit{Prefix raises $p$(\texttt{no}) and gives a correct top-1 prediction.}
\end{quote}

\textbf{Causal Direction (\texttt{Does \{cause\} lead to \{effect\}? Answer yes or no:})}

Causal direction shows the strongest response to soft prompting. Accuracy increases
from 0\% to 53.8\% (7 of 13 prompts) and mean $p(\text{correct})$ increases by $+0.236$.
As with balanced parentheses, the baseline model fails at output format, predicting
\texttt{`1'} as the top token on every prompt. The prefix resolves this entirely,
shifting predictions to \texttt{`Yes'} or \texttt{`No'} and achieving above-chance
accuracy on the semantic direction judgement itself.
The concept involves asymmetric causal plausibility
(e.g.\ ``warmth $\to$ light'' is implausible while ``storm $\to$ flood'' is
plausible), and the prefix correctly classifies 7 of 13 such pairs despite
consisting of only 10 learned tokens with no base-model parameter updates.

\begin{quote}
\fontfamily{pcr}\selectfont\small
\textbf{Prompt:} \textup{Does warmth lead to light? Answer yes or no: }\hfill[\textit{correct: no}]\\
\textbf{Baseline:}\ `1' (0.67) \quad `Yes' (0.18) \quad `No' (0.05)\\
\textbf{Prefix:}\ \ \ `No' (0.55) \quad `Yes' (0.20) \quad \ldots\\
\textit{Prefix corrects format and semantic direction ($p$ rises from 0.051 to 0.578).}\\[6pt]
%
\textbf{Prompt:} \textup{Does storm lead to flood? Answer yes or no: }\hfill[\textit{correct: yes}]\\
\textbf{Baseline:}\ `1' (0.48) \quad `Yes' (0.18) \quad `No' (0.16)\\
\textbf{Prefix:}\ \ \ `Yes' (0.58) \quad `No' (0.20) \quad \ldots\\
\textit{Prefix correctly identifies causal direction with high confidence.}\\[6pt]
%
\textbf{Prompt:} \textup{Does frost lead to death? Answer yes or no: }\hfill[\textit{correct: yes}]\\
\textbf{Baseline:}\ `1' (0.36) \quad `Yes' (0.16) \quad `No' (0.12)\\
\textbf{Prefix:}\ \ \ `No' (0.37) \quad `Yes' (0.11) \quad \ldots\\
\textit{Prefix fixes format but predicts the wrong direction; a failure case.}
\end{quote}

\subsection{Discussion}

The three concepts reveal a consistent pattern. For balanced parentheses and causal
direction, the baseline model fails entirely at producing the required output
format: the top-1 token is \texttt{`1'} regardless of the input, indicating that
the model does not interpret the yes/no instruction from the template alone.
The soft prefix resolves the format failure in both cases, and for causal direction
it additionally steers the semantic direction of the response above chance. The
relatively large gain on causal direction ($+53.8\%$ accuracy, $+0.236$ in
$\bar{p}$) compared with balanced parentheses ($+28.6\%$, $+0.078$) may reflect
that causal plausibility judgements are more accessible to a continuous prefix
operating on early residual-stream states, whereas recursive bracket matching
requires tracking a global structural property that a fixed 10-token prefix cannot
fully supply.

Residue class stands apart. The baseline already produces numeric outputs, since
the template ends with \texttt{=} and thereby elicits digit completions, and the
concept is already evaluated at non-trivial accuracy (15\%). The prefix fails to
improve it, consistent with the absence of a localised binary concept direction
in the residual stream analysis of Section~\ref{sec:residue_results}.
These three cases together suggest that soft prompting is most effective when the
base model's failure is one of output framing rather than of knowledge. The prefix
can redirect the model to the correct response format and, when the underlying
representation is accessible, to the correct answer, but it cannot supply knowledge
that is absent or structurally inaccessible at the residual stream level.

% \section{...}

% % Maybe mention that/why stitching doesn't work, or GRU steering (run experiment again)

